% ===== PRÉAMBULE CANONIQUE — à coller VERBATIM en tête de CHAQUE .tex =====
\documentclass[11pt,a4paper]{article}
\usepackage[utf8]{inputenc}\usepackage[T1]{fontenc}\usepackage{lmodern}
\usepackage[french]{babel}
\usepackage{amsmath,amssymb,mathtools}\usepackage{icomma}
\usepackage[version=4]{mhchem}          % \ce{...} : équations & formules
\usepackage{siunitx}\sisetup{locale=FR} % \qty{}{}, \si{}, \ang{}
\usepackage{chemfig}                    % \chemfig{...} : structures développées
\usepackage{xcolor}\usepackage{colortbl}
\usepackage{tikz}\usepackage{pgfplots}\pgfplotsset{compat=1.18}
\usepackage[most]{tcolorbox}\usepackage{enumitem}\usepackage{array,booktabs}
\usepackage[a4paper,margin=2.2cm]{geometry}
\usetikzlibrary{arrows.meta,calc,angles,quotes,positioning,patterns,backgrounds,shapes.geometric}
\definecolor{primary}{RGB}{222,49,99}\definecolor{primarydark}{RGB}{150,22,55}
\definecolor{defcol}{RGB}{0,114,178}\definecolor{methcol}{RGB}{52,92,148}
\definecolor{excol}{RGB}{0,135,90}\definecolor{attncol}{RGB}{200,40,55}
\tcbset{boxbase/.style={enhanced,breakable,boxrule=0.9pt,arc=2.5pt,
  left=3mm,right=3mm,top=2.3mm,bottom=2.3mm,fonttitle=\bfseries,coltitle=white}}
% --- boîtes TUTORIEL (noms EXACTS, ne pas renommer) ---
\newtcolorbox{cle}[1][]{boxbase,colback=defcol!6,colframe=defcol,
  title={\textsc{À retenir}\ifx\relax#1\relax\relax\else\textmd{ : \itshape #1}\fi}}
\newtcolorbox{methode}[1][]{boxbase,colback=methcol!7,colframe=methcol,sharp corners,
  title={\textsc{Méthode}\ifx\relax#1\relax\relax\else\textmd{ --- \itshape #1}\fi}}
\newtcolorbox{exemplebox}[1][]{boxbase,colback=excol!5,colframe=excol,
  title={\textsc{Exemple}\ifx\relax#1\relax\relax\else\textmd{ : \itshape #1}\fi}}
\newtcolorbox{piege}[1][]{boxbase,colback=attncol!6,colframe=attncol,
  title={\textsc{Piège}\ifx\relax#1\relax\relax\else\textmd{ : \itshape #1}\fi}}
\newtcolorbox{remarque}[1][]{boxbase,colback=black!4,colframe=black!45,
  title={\textsc{Remarque}\ifx\relax#1\relax\relax\else\textmd{ : \itshape #1}\fi}}
% --- boîtes EXERCICE / CORRIGÉ (noms EXACTS, auto-comptées) ---
\newtcolorbox[auto counter]{exo}[1][]{boxbase,colback=primary!4,colframe=primary,
  title={\textsc{Exercice}~\thetcbcounter\ifx\relax#1\relax\relax\else\textmd{ --- \itshape #1}\fi}}
\newtcolorbox[auto counter]{corrige}[1][]{boxbase,colback=excol!4,colframe=excol,
  title={\textsc{Corrigé}~\thetcbcounter\ifx\relax#1\relax\relax\else\textmd{ --- \itshape #1}\fi}}
\newcommand{\R}{\mathbb{R}}\newcommand{\N}{\mathbb{N}}\newcommand{\Z}{\mathbb{Z}}\newcommand{\ds}{\displaystyle}
% (un \usepackage{fancyhdr}+en-tête est possible ; il est ignoré côté web)
% =========================================================================
\begin{document}
% THEME_TITLE: Piles galvaniques et électrolyse
\sisetup{per-mode=symbol}
\begin{center}\color{primarydark}\Large\bfseries Thème 5 — Piles galvaniques et électrolyse\end{center}

Une \textbf{pile} convertit l'énergie d'une réaction redox \textbf{spontanée} en énergie électrique.
Les deux demi-réactions sont \textbf{séparées} : les électrons doivent passer par un circuit
extérieur.

\begin{cle}[anode et cathode d'une pile]
\begin{itemize}[leftmargin=1.5em,topsep=1pt,itemsep=0pt]
  \item \textbf{Anode} (pôle $-$) : siège de l'\textbf{oxydation}, elle \emph{produit} les électrons.
        C'est le couple de $E_{\mathrm{r}}^{\circ}$ le \textbf{plus bas}.
  \item \textbf{Cathode} (pôle $+$) : siège de la \textbf{réduction}, elle \emph{consomme} les
        électrons. C'est le couple de $E_{\mathrm{r}}^{\circ}$ le \textbf{plus haut}.
\end{itemize}
La pile débite si $\Delta E^{\circ}=E_{\mathrm{r}}^{\circ}(\text{cathode})-E_{\mathrm{r}}^{\circ}
(\text{anode})>0$.
\end{cle}

\begin{center}
\begin{tikzpicture}[scale=0.82,>=Stealth,every node/.style={transform shape}]
  % béchers
  \draw[gray,line width=1pt] (-4,0.7)--(-4,-1.7)--(-2,-1.7)--(-2,0.7);
  \fill[blue!11] (-4,-1.68) rectangle (-2,0.1);
  \draw[gray,line width=1pt] (2,0.7)--(2,-1.7)--(4,-1.7)--(4,0.7);
  \fill[green!11] (2,-1.68) rectangle (4,0.1);
  % électrodes
  \fill[violet!65] (-3.1,1.9) rectangle (-2.9,-1.3);
  \node[white,font=\footnotesize\bfseries,rotate=90] at (-3,0.3){Zn};
  \fill[orange!80] (2.9,1.9) rectangle (3.1,-1.3);
  \node[white,font=\footnotesize\bfseries,rotate=90] at (3,0.3){Cu};
  \node[blue!55!black,font=\scriptsize] at (-3,-1.05){\ce{Zn^{2+}}};
  \node[green!45!black,font=\scriptsize] at (3,-1.05){\ce{Cu^{2+}}};
  % fils + voltmètre
  \draw[gray,line width=1pt] (-3,1.9)--(-3,3.0)--(-0.8,3.0);
  \draw[gray,line width=1pt] (3,1.9)--(3,3.0)--(0.8,3.0);
  \draw[gray,line width=1pt,fill=white] (-0.8,2.6) rectangle (0.8,3.4);
  \node[font=\bfseries,primarydark] at (0,3.0){V};
  % flux d'électrons (anode -> cathode dans le fil)
  \draw[->,blue,line width=1.1pt] (-2.5,3.22)--(-1.4,3.22);
  \node[blue,font=\scriptsize] at (-1.95,3.5){\ce{e^-}};
  \draw[->,red!70,line width=0.9pt,dashed] (1.4,2.78)--(2.5,2.78);
  \node[red!70,font=\scriptsize] at (1.95,2.55){$i$};
  % pont salin
  \draw[orange!55,line width=5pt,line cap=round] (-2.35,-0.2) .. controls (-2.35,1.5) and (2.35,1.5) .. (2.35,-0.2);
  \node[orange!45!black,font=\scriptsize] at (0,1.55){pont salin};
  \draw[->,gray,line width=0.7pt] (-0.5,1.15)--(-1.1,1.15);
  \node[gray,font=\tiny] at (-0.8,1.38){anions};
  \draw[->,gray,line width=0.7pt] (0.5,1.15)--(1.1,1.15);
  \node[gray,font=\tiny] at (0.8,1.38){cations};
  % légendes demi-équations
  \node[violet!55!black,font=\scriptsize,align=center] at (-3,-2.35){\textbf{Anode} ($-$)\\ \ce{Zn -> Zn^{2+} + 2 e^-}};
  \node[orange!55!black,font=\scriptsize,align=center] at (3,-2.35){\textbf{Cathode} ($+$)\\ \ce{Cu^{2+} + 2 e^- -> Cu}};
\end{tikzpicture}
\end{center}

\begin{cle}[sens des porteurs et notation conventionnelle]
Dans le \textbf{fil}, les \textbf{électrons} vont de l'\textbf{anode ($-$) vers la cathode ($+$)} ; le
courant $i$ circule en sens inverse. Dans le \textbf{pont salin}, les \textbf{anions} migrent vers
l'anode et les \textbf{cations} vers la cathode (électroneutralité). On écrit la pile :
\[
(\,-\,)\ \text{Anode}\ \big|\ \text{solution}\ \big\|\ \text{solution}\ \big|\ \text{Cathode}\ (\,+\,)
\qquad\text{ex. }\ \ce{(-) Zn | Zn^{2+} || Cu^{2+} | Cu (+)}
\]
($|$ = interface électrode/solution, $\big\|$ = pont salin.)
\end{cle}

\begin{cle}[électrolyse : forcer une réaction non spontanée]
Dans une \textbf{électrolyse}, un \textbf{générateur} impose le courant pour réaliser une réaction
\textbf{non spontanée}. La cathode est reliée au pôle $-$ du générateur, l'anode au pôle $+$.
\textbf{Règle d'or} (pile \emph{comme} électrolyse) : la \textbf{réduction} a toujours lieu à la
\textbf{cathode}, l'\textbf{oxydation} à l'\textbf{anode} ; seul le \emph{signe} des électrodes change.
\end{cle}

\begin{cle}[loi de Faraday]
La masse déposée (ou transformée) à une électrode vaut
\[
\boxed{\,m=\dfrac{M\,I\,t}{n\,F}\,}\qquad F=\qty{96485}{C\per mol},\quad Q=I\,t,
\]
avec $m$ en g, $M$ en g/mol, $I$ en \textbf{ampères}, $t$ en \textbf{secondes}, $n$ le nombre
d'électrons de la demi-équation. ($Q=It$ est la charge, en coulombs ; $Q/F$ = moles d'électrons.)
\end{cle}

\begin{exemplebox}[dépôt de cuivre]
Électrolyse de \ce{CuSO4}, $I=\qty{2.0}{A}$, $t=\qty{30}{min}=\qty{1800}{s}$. À la cathode
\ce{Cu^{2+} + 2 e^- -> Cu}, donc $n=2$ et $M(\ce{Cu})=\qty{63.5}{g\per mol}$ :
\[
m=\frac{63{,}5\times 2{,}0\times 1800}{2\times 96485}\approx \qty{1.18}{g}.
\]
\end{exemplebox}

\begin{piege}[erreurs fréquentes]
\begin{itemize}[leftmargin=1.5em,topsep=1pt,itemsep=0pt]
  \item Convertir $t$ en \textbf{secondes} (min $\times 60$, h $\times 3600$) avant d'appliquer Faraday.
  \item Les \textbf{électrons ne traversent jamais la solution} : dans la solution et le pont salin,
        ce sont les \textbf{ions} qui transportent le courant.
  \item $n$ est le nombre d'électrons de la demi-équation ($n=2$ pour \ce{Cu^{2+}}, $n=1$ pour \ce{Ag+}).
\end{itemize}
\end{piege}

\end{document}
